\cotefile{161-4}
\scanpage{18}

donc la question est si c'est une équiv. de topos

\[
\widehat{C\times D}\overset{?}{\simeq}\widehat C\times\widehat D \tag*{(2)}
\]

Si $C$ et $D$ sont stables par $\varprojlim$ finies, alors $C\times D$ l'est aussi. \deleted{\unsure{Or si $D$|et soit $A$}} de façon générale $A$ est une telle catégorie, on sait que

\[
\begin{aligned}
\underline{\mathrm{Hom}}_{\mathrm{top}}(E,A)&\simeq
  (\underline{\mathrm{Hom}}^{\unsure{g|?|*}}_{\unsure{\mathrm{gex}|\mathrm{sex}}}(A,E))^{\mathrm o}\\
u&\longmapsto u^*
\end{aligned}
\]

\deleted{et} \inserted{on conclut} \uline{dans ce cas}, (2) est une équivalence.

c) Dans le cas général, on peut écrire pour des morphismes

\[
\deleted{\unsure{E}\quad i:C\longrightarrow D\qquad j:E\longrightarrow\widehat D}
\]

\[
i:C\longrightarrow C'\qquad j:E\longrightarrow\widehat D
\qquad\text{(i.e. $E\simeq\widetilde D$ pour une top. conv. sur $D$)}
\]

$i$ pl. fidèle, $j$ plongement de topos, \inserted{$C'$ et $D$ stbles par $\varprojlim$ finies,} \deleted{et \unsure{on|on en}} on déduit un diagramme commutatif mod iso. can.

\diagram{\begin{tikzcd}[ampersand replacement=\&,column sep=huge,row sep=large]
\pi(C,\widetilde D) \arrow[r,"{\alpha=\pi(i,j)}"] \arrow[d,"\gamma"'] \&
\pi(C',\widehat D)\simeq\widehat{C'\times D}\arrow[d,"{\simeq\ \gamma'}"]\\
\widehat C\times\widetilde D\arrow[r,hook,"{\beta=i\times j}","{\substack{\text{plongement}\\\text{de topos comme}\\\text{le pr. direct de tels}}}"'] \& \widehat{C'}\times\widehat D
\end{tikzcd}}

\margin{Si $\gamma$ est une équivalence, alors $\alpha$ est un plongement. Peut-on vérifier que $\alpha$ l'est en effet ? \deleted{\unsure{Oui, car}}}

Or $\alpha=\pi(i,j)$ est un composé

\diagram{\begin{tikzcd}[ampersand replacement=\&,column sep=large]
\pi(C,\widetilde D)\arrow[r,"\alpha_0"]\&\pi(C,\widehat D)\arrow[r,"\alpha_1"]\arrow[d,phantom,"\rotatebox{90}{$\simeq$}"]\&\pi(C',\widehat D)\arrow[d,phantom,"\rotatebox{90}{$\simeq$}"]\\
\&\widehat{C\times D}\arrow[r,"{\substack{\text{plongement}\\\text{car } C\times D\to C'\times D\ \text{pl. fid.}}}"']\&\widehat{C'\times D}
\end{tikzcd}}

donc on est ramené à voir seulement si

\diagram{\begin{tikzcd}[ampersand replacement=\&,row sep=tiny]
\alpha_0:\pi(C,\widetilde D)\arrow[r]\&\pi(C,\widehat D)\arrow[d,phantom,"\rotatebox{90}{$\simeq$}"]\&\text{un plongement}\\
\&\widehat{C\times D}\&
\end{tikzcd}}

si le foncteur image \deleted{inverse associé est} directe associé est

\[
\begin{aligned}
&\deleted{F\longmapsto (X\longmapsto \unsure{a|\underline a}(FX))}\\
&u\longmapsto j_*\circ u\qquad\text{qui est bien pl. fidèle}
\end{aligned}
\]

\uline{Corollaire} \deleted{\unsure{Il résulte}}, $\gamma(C,E):\pi(C,E)\longrightarrow\widehat C\times\widetilde D$ est un \uline{plongement}.

(Car $\gamma'\alpha$ l'est, et $\beta$ aussi, donc $\beta\gamma$ l'étant, $\gamma$ aussi)

Notons que \deleted{\unsure{l'est}} $\widehat C\times\widetilde D$ est la sous-catégorie \inserted{des} pleine de $\widehat {C'}\times\widehat D\simeq\widehat{C'\times D}$, \deleted{\unsure{décrite}} alors formée des $F\in\mathrm{Ob}\,\widehat{C'\times D}$ tels que (désignant par $p$ et $q$ les deux projections $\widehat{C'\times D}\ \substack{\to\widehat{C'}\\[-2pt]\to\widehat D}$) on ait

\[
\left\{\begin{aligned}
p_*(F)&\in\widehat C\ \text{(identifié à la sous-catégorie strict. pleine image essentielle de $i_*:\widehat C\to\widehat{C'}$)}\\
q_*(F)&\in\widetilde D
\end{aligned}\right.
\]

Prouvons que c'est \deleted{\unsure{tel}} (l'image ess. de) $\pi(C,D)$, ce qui prouvera que $\gamma$ est une équivalence.

% archive sheet number: 17
\scanpage{20}

\uline{Géom. Algébrique}

Compléments sur la formation des $\varprojlim$ et des sommes dans $\mathrm{Aff}_k$.

a) \uline{Les sommes \inserted{finies} sont disjointes et universelles}

\[
X=\coprod_{i\in I}X_i\longleftarrow X_i
\]

1) $X_i\to X$ un mono et $=$ une \uline{imm. fermée} \qquad ($A=\prod A_i\to A_i$ surjectif)

2) $X_i\times_X X_j=\varnothing$ si $i\ne j$ \qquad i.e. $A_i\otimes_A A_j=0$ si $i\ne j$ car l'\unsure{idemt|élément|égal} $1$ de $A_i$, $0$ de $A_j$ \deleted{\unsure{égal à $1$ dans $A_j$, $0$ dans $A_i$}}

\diagram{\begin{tikzcd}[ampersand replacement=\&,column sep=normal,row sep=small]
3)\quad\forall\quad X'\arrow[r]\&X\&\text{on a}\quad X'\xleftarrow{\sim}\coprod X'_i\\
X'_i\arrow[u]\arrow[r]\&X_i\arrow[u]\&\unsure{q.}
\end{tikzcd}}

i.e. $\forall\,A=\prod\unsure{A_i|A_j}\to B$, $B$ décomposé en $\prod B_i$, et \ldots

\deleted{\unsure{Donc dans ce cas|b) Dans le cas où}} Donc les produits \inserted{(ou \uline{produits fibrés})} sont distributifs p. r. aux sommes

\[
\left(\coprod X_i\right)\times\left(\coprod_j Y_j\right)\simeq\coprod_{ij}(X_i\times Y_j)
\]

\[
\bigl[\,\deleted{\unsure{Y}}X\times\underbrace{\coprod Y_i}_{Y}\longleftarrow\coprod_i X\times Y_i\qquad\ldots\,\bigr]
\]

\diagram{\begin{tikzcd}[ampersand replacement=\&,row sep=small]
X\times Y_i\arrow[r]\arrow[d]\&X\times Y\arrow[d]\&X\arrow[d,dash]\\
Y_i\arrow[r]\&Y\arrow[r,dash]\&e
\end{tikzcd}}

Donc le foncteur

\[
\begin{aligned}
(\mathrm{Ens}\,\unsure{\mathrm{fin}|\mathrm f|\mathrm{finis}})&\longrightarrow\mathrm{Aff}_k\\
\underline I&\longrightarrow I_k=\coprod_I e_k
\end{aligned}
\]

commute aux \deleted{$\unsure{\text{limites}}\ \varprojlim$ finies} produits \inserted{finis}, \deleted{\unsure{en particulier aux}} \inserted{et $=$ aux} produits fibrés, donc aux $\varprojlim$ finies. Donc transforme monoïde (groupe) fini en monoïde (groupe) algébrique affine.

Du point de vue alg. comm., on a

\margin{Le foncteur $I_k$ est le foncteur : ens des décompositions idempotentes indexées par $I$}

\[
I_k=V_{k^I}\qquad k\to k^I\quad\text{contravariant en $I$}
\]

\inserted{algèbre produit $\simeq\mathrm{Fonc}(I,k)$}

$k^I\simeq k[I]$ \inserted{(pr. finis)} (gare aux confusions avec polynômes, Si \deleted{\unsure{$I$ a}} \inserted{$I$ a} une \deleted{\unsure{grpe, la multiplication}} \inserted{loi de composition} interne $I\times I\to I$, cela définit

\[
k[I]\longrightarrow k[I\times I]\simeq k[I]\otimes k[I]
\]

la comultiplication \ldots.

% archive sheet number: 19
\scanpage{21}

\[
\mathrm{Spec}\,I_k=\mathrm{Spec}\,k^I=\coprod_I\underbrace{\mathrm{Spec}(k)}_{S}
\qquad\text{\unsure{envisagé|envoyé} dans \unsure{Sch|$S$} de façon évidente.}
\]

b) Produits fibrés dans le cas où l'un des deux morphismes est un monomorphisme. Alors le diagramme correspondant sur les spectres est cartésien

\diagram{\begin{tikzcd}[ampersand replacement=\&,row sep=small]
\mathscr X\arrow[r,dashed,leftrightarrow]\arrow[d]\&\mathscr X'\arrow[d,dotted]\\
\mathscr Y\arrow[r,leftrightarrow]\&\mathscr Y'
\end{tikzcd}}

\diagram{\begin{tikzcd}[ampersand replacement=\&,row sep=small]
X\arrow[r]\arrow[d]\&X'\arrow[d]\\
Y\arrow[r]\&Y'
\end{tikzcd}}

De plus, si \uline{$Y'\to Y$} est une imm. fermée (resp. ouverte) -- plus généralement si cela correspond à $A\to AS^{-1}$, donc à un homéomorphisme dans, alors \deleted{le diagramme} de $=$ pour $\mathscr X'\unsure{\leftarrow|-|\to}\mathscr X$, donc la top. de $X'$ est la top. induite.

\uline{Cas particulier} Intersection de deux sous-objets ---

\diagram{\begin{tikzcd}[ampersand replacement=\&,row sep=small]
c)\quad A\arrow[r,"u"]\& B\\
Y\&X\arrow[l,"\varphi"]
\end{tikzcd}}

formule pour $\overline{\varphi(Z)}$, $Z=V(J)$

Condition pour que $\varphi$ soit dominant :

\[
\mathrm{Ker}\,\varphi\subset\mathrm{Nil}\,A
\]

% archive sheet number: 20
\scanpage{22}

\uline{Topologie des spectres}

Remarque : espace noethérien $=$ \deleted{tout} espace quasicompact tel que a) tout sous-espace fermé ait un nb fini de comp. irréductibles b) toute suite décroissante de sous-espaces fermés \unsure{irred} est stationnaire.

Composantes connexes d'un espace top. n'ayant qu'un nb fini de composantes irréductibles (p. ex. noeth.) -- plus gén. où l'ens des comp. irréd. est loc. fini.

\deleted{Topologie de $\mathrm{Spec}\,A_f$ \unsure{et $X_f$} L'espace topologique $\mathrm{Spec}\,A$ pour $A$ variable.} base de topologie formée des $X_f$ \deleted{/ homéomorphismes $X_f\simeq\mathrm{Spec}\,A_f$}

L'espace topologique $\mathrm{Spec}\,A$ pour $A$ variable \quad Cf.

\diagram{\begin{tikzcd}[ampersand replacement=\&,row sep=small]
A\arrow[r]\&B\\
Y=\mathrm{Sp}\,A\&X=\mathrm{Spec}\,B\arrow[l,"\varphi"]
\end{tikzcd}}

\[
\left\{\begin{aligned}
\varphi(x_{\mathfrak p})&=y_{\varphi^{-1}(\mathfrak p)}\\
\varphi^{-1}(V_{\deleted{\unsure{A}}}(J))&=V(u(J)B)\\
\varphi^{-1}(Y_f)&=X_{\deleted{\varphi}u(f)}
\end{aligned}\right.
\]

\uline{Ex} a) \deleted{\unsure{f}} $u:A\to B$ surjectif, $\varphi$ imm. fermée

b) \deleted{\unsure{f}} $u:A\to A_f$ \qquad $\varphi$ immersion ouverte \qquad $\mathrm{Spec}(A_f)\simeq X_f$

dans les deux cas, les ext. résiduelles sont triviales, car $A\to B$ un épim. d'anneaux [i.e. $\mathscr X\to\mathscr Y$ un monomorphisme d'espaces alg. affines]

\[
\begin{aligned}
\text{c)}\quad \mathrm{Spec}\,\prod_I A_i&\simeq\coprod_i\mathrm{Spec}(A_i)\\
X&\simeq\coprod_i X_i
\end{aligned}
\]

d) \uline{Attention, le foncteur} $\mathscr X\to\mathrm{Spec}\,\mathscr X$ \uline{ne commute pas aux produits ou produits fibrés} [\unsure{même|vrai} sur un corps alg. clos]

% slanted
\margin{donner exemple du produit \unsure{$E'\times E'$|$kE'\times E'$|$\mathbb E^1\times\mathbb E^1$}}

\diagram{\begin{tikzcd}[ampersand replacement=\&,row sep=small,column sep=small]
\&\&z=V_{K'}\arrow[dl,dash]\arrow[dr,dash]\&\&\\
V_{k(x)}\arrow[r,dash]\arrow[d]\&V_A\&\&V_B\&V_{k(y)}\arrow[l]\arrow[d]\\
x\arrow[r]\&X\arrow[dr]\&\&Y\arrow[dl,dash]\&y\arrow[l]\\
\&\&\substack{S\\ V_k}\&\&
\end{tikzcd}}

\diagram{\begin{tikzcd}[ampersand replacement=\&,row sep=tiny,column sep=tiny]
\& C\arrow[dl,dash]\arrow[dr,dash]\&\\
A\arrow[dr,dash]\&\&B\arrow[dl,dash]\\
\&k\&
\end{tikzcd}}

\diagram{\begin{tikzcd}[ampersand replacement=\&,row sep=tiny,column sep=tiny]
\&\mathscr X\times_S\mathscr Y\arrow[dl]\arrow[dr]\&\\
x\in\mathscr X\arrow[dr,dash]\&\&\mathscr Y\ni y\arrow[dl,dash]\\
\&s\in S\&
\end{tikzcd}}

\[
\left|\mathrm{Spec}\,\mathscr X\times_S\mathscr Y\right|\longrightarrow
\left|\mathrm{Spec}(\mathscr X)\right|\times_{|\mathrm{Spec}\,S|}\left|\mathrm{Spec}\,\mathscr Y\right|
\]

\diagram{\begin{tikzcd}[ampersand replacement=\&,row sep=tiny,column sep=tiny]
\&K'\ \text{corps}\arrow[dl,dash]\arrow[dr,dash]\&\\
k(x)\arrow[dr,dash]\&\&k(y)\arrow[dl,dash]\\
\&k(s)\&\\
\&k\arrow[u]\&
\end{tikzcd}}

cette application est surjective, mais n'est pas injective. la \deleted{\unsure{image de}} fibre de $x,y$ \deleted{\unsure{sont}} est can. équivalente à $x\times_S y\simeq\mathrm{Spec}\,k(x)\otimes_{k(s)} k(y)$ [qui est $\ne\varnothing$]

\deleted{Mais} \inserted{Pour esp. alg. affines de t.f.,} sur un corps alg clos, la \inserted{formation du} Spec Maximal \inserted{ensembliste} commute aux $\varprojlim$ finies, mais la top. n'est pas la topologie produit.

% archive sheet number: 21
\scanpage{23}

Prerequisites for the courses I and II

I) Introduction to A.G. : Category theory, basic algebra including some commutative algebra. Bourbaki is a good reference. Most facts we are using will be recalled.

II) Introduction to algebraic groups. Familiarity with the basic notions of A.G. will be assumed.

A.G. is a messy subject, \inserted{\unsure{even upon foundations|even more foundations},} we cannot give anything like a detailed exposition of all of it in one summer course. We will have to give rather a survey of some key points, being rather sketchy for proofs, -- most of which are straightforward. We can \unsure{pile in|pile up|pick up|fill in} whatever is requested. Homework will be necessary for those who do not have already a fair background.

Question is how to \unsure{relate|belt} the two courses. The contents of \deleted{\unsure{II}} \unsure{I are|are} prerequisites for II !

\textcircled{1}\quad $k$

\[
P_I=P[(T_i)_{i\in I}]
\]

\[
k^I\simeq\mathrm{Hom}_{k\text{-}\mathrm{alg}}(P_I,k)
\qquad\inserted{\text{l'espace affine sur $k$ de dim $I$}}\quad\deleted{\unsure{V_{\mathrm{alg}}(P_I,k)}}
\]

\margin{$I,J$ finis ou infinis}

\[
\begin{aligned}
&f\in P_I,\quad x=(a_i)_{i\in I}\in k^I\qquad\deleted{\unsure{f\in k^{\prime I}}}\\
&f(x)=\varepsilon_x(f)
\end{aligned}
\]

\[
F_j(t)=0\qquad j\in J\qquad V((F_j)) \tag*{(1)}
\]

\textcircled{2} (a) Pbs \inserted{du type} arithmétiques : très \unsure{différents|affines|différent} suivant la nature de $k$ : \ldots, $\mathbb R,\mathbb C$

\textcircled{3} (b) Pbs \inserted{de type} géométriques, issu de la géométrie sur $\mathbb R$ \deleted{et} sur $\mathbb C$ on cherche les solutions de (1) \deleted{pour} \uline{toute} \deleted{extension} \inserted{$k$-algèbre} $k'$, \deleted{\unsure{de $k$}} et comment elles varient ; et les propriétés de (1) qui sont stables par de tels changements \ldots

\[
\begin{aligned}
\mathbb E_k^I:\quad k'&\longmapsto k^{\prime I} &\mathbb E_k^I(k')&=k^{\prime I}\\
V_S(k')&\subset\mathbb E_k^I(k') & V_S&\subset\mathbb E_k^I
\left\{\begin{aligned}&\text{étude indépendante du plongement}\\ &\text{étude \unsure{du plongement|des plongements}}\end{aligned}\right.
\end{aligned}
\]

% archive sheet number: 22
\scanpage{24}

\uline{Sous-espace \deleted{\unsure{algébr.|algébrique}} fermé de $\mathbb E_k^I$} \uline{espace alg. affine}

\[
\text{\textcircled{4}}\quad V_S=V_J
\]

\inserted{idéal engendré, c'est le syst. d'équations le plus naturel}

\margin{référence au ``th. de la base'' de Hilbert ($k$ corps \deleted{\unsure{plus}} gén, noeth) $\Rightarrow k[T_1-\ T_n]$ noeth}

\[
A_J=P_I/J
\]

\diagram{\begin{tikzcd}[ampersand replacement=\&,row sep=small,column sep=small]
k^{\prime I}\arrow[r,phantom,"\simeq"]\arrow[d,phantom,"\rotatebox{90}{$\subset$}"]\&\mathrm{Hom}_{k\text{-}\mathrm{alg}}(P_I,k')\arrow[d,phantom,"\rotatebox{90}{$\subset$}"]\\
V_J(k')\arrow[r,phantom,"\simeq"]\&\{\deleted{\unsure{\mathrm{hom}}}u:P_I\to k'\mid\deleted{\unsure{u\ \text{nul sur }J}}\inserted{u(J)=0}\ (\inserted{\text{$u$ nul}}\ \text{sur les }F_j)\}\arrow[d,phantom,"\rotatebox{90}{$\simeq$}"]\\
\&\mathrm{Hom}_k(A_J,k')
\end{tikzcd}}

\deleted{\unsure{4}} Les foncteurs de la forme $V_S$ ($I,S$ convenables) sont, à isom près, les foncteurs représentables $V_A$ donner un isom. entre $V_A$ et un $V_S$ \deleted{\unsure{rev}} revient au $=$ que de se donner un isom $A\simeq P_I/J$ ($J$ idéal de $P_I$ engendré par $S$), \deleted{\unsure{donc}} ; se donner un plongement $V_A\subset\mathbb E_k^I$ \deleted{\unsure{revient à se donner un}} \inserted{identifiant $V_A$ à un sous-esp. alg fermé de $\mathbb E_k^I$}

% slanted
\margin{utilisant le \unsure{résultat sur les|résult. sur les|notion de} foncteurs représentables}

[Plus généralement, tout plongement $V_A\hookrightarrow\deleted{\unsure{E}}\mathbb E^I$ correspond à un épimorphisme d'anneaux $P_I\to A$ (mais un tel épimorphisme n'est pas en général surjectif \ldots] s'identifient aux hom \deleted{\unsure{surjectifs}} $P_I\to A$, ou encore aux systèmes de gén. $(a_i)_{i\in I}$ de $A$ comme \uline{algèbre} sur $k$.

\margin{Rem}

\textcircled{5} \uline{Cor} La cat. des esp. alg. affines sur $k$ est anti-équivalente à celle des algèbres sur $A$. Ceux qui se plongent comme sous-espaces alg. fermés dans quelque $\mathbb E_k^n$ ($n$ entier $\geq0$ fini) \deleted{sont ceux} correspondent aux $A$ qui \deleted{\unsure{sont un syst. de générateurs finis|sont des $k$-alg. de type fini}} \inserted{sont des $k$-alg. de type fini}.

Les \deleted{\unsure{idéaux de}} sous-espaces algébriques de $\mathbb E_k^I$ correspondent aux idéaux de $P_I$. \deleted{\unsure{en}} $\bigl[$ \uline{\unsure{si on connait|si on se donne|$J$ connu quand on} $=$ connait} $V_J$ \ldots \deleted{\unsure{$J'$}}$J$

\[
\begin{aligned}
V_J\subset V_{J'}&\Longleftrightarrow J'\subset J\\[-2pt]
&\Leftarrow\left\{\begin{aligned}&=\{f\in P_I\mid f(x)=0\ \ \forall\\&\qquad k'/k\text{ et }x\in V_J(k')\}\end{aligned}\right.
\end{aligned}
\]

% archive sheet number: 23
\scanpage{25}
% archive sheet number: 24

\textcircled{6} Ce qui se passe si on se restreint au cas où $k'$ est un corps [on est tenté de le faire au moins si $k$ en est un]. \deleted{\unsure{la correspondance}} Soit $V'_S=V_S/\text{corps sur }k$

\[
V^{\prime}_{\deleted{\mathrm{corps}},S}=V^{\prime}_{\deleted{\mathrm{corps}},J}=V^{\prime}_{\deleted{\mathrm{corps}},\widetilde J}
\]

Donc $J\mapsto V^{\prime}_{\deleted{\unsure{\mathrm{corps}}}J}$ plus injective. Mais sa restriction à l'ens des idéaux \deleted{\unsure{J}} $J$ tels que $J=\widetilde J$ l'est, car on a

\[
\widetilde J=\{f\in P_I\mid f(x)=0\quad\forall\,x\in V'_J(k'),\ k'/k\text{ un corps}\}
\]

[Cela \deleted{signifie} équivaut à l'énoncé suivant d'AC

\uline{Prop} $(\forall\,\underset{\text{idéal}}J\subset P_I)$ $\widetilde J$ est l'intersection des noyaux des hom de $P_I/J$ dans des corps

\[
\Downarrow
\]

$\forall$ anneau $A$, l'intersection des hom de $A$ dans des corps [i.e. l'intersection des idéaux premiers) est l'ens des éléments nilpotents. ] Donc

\[
\widetilde J\subset\widetilde J'\ \unsure{\longleftrightarrow|\Longleftrightarrow}\ V(J')\unsure{\subset|=}\,V(J)
\]

\unsure{M$^{\mathrm s}$|Mêmes} conclusions \inserted{pour $I$ fini} si on se borne à tte sous-catégorie \inserted{$C$} de la catégorie des $k$-algèbres réduites, pourvu toutefois que \unsure{$C$|$\mathcal C$|$k$} soit assez grande pour que l'analogue de la prop. précédente soit valable

\textcircled{$6'$} (par ex. $C$ contient les extensions de corps,] -- \deleted{\unsure{ou si}} si $I$ est fini, s'il contient les ext. finies des corps résiduels de $k$, -- p. ex. ou \unsure{aust. : $k$ si|aust : $k$ si|aura : $k$ si} $k$ est un corps alg. clos \ldots)

\margin{\uline{Nullstellensatz}}

\margin{Dans une alg. de t.f. sur corps $k$ tt idéal premier (réduit) est intersection d'idéaux maximaux, et le quot\unsure{ient} par un \ldots{} est \ldots}

\textcircled{7} Importance des \deleted{\unsure{points}} éléments nilpotents.

\[
\text{\textcircled{8}}\quad A\simeq\mathrm{Hom}_{k\text{-}\mathrm{alg}}(k[T],A)\simeq\mathrm{Hom}(V_A,\mathbb E_k^1)
\]

% slanted
\margin{définir d'abord $f(x)$ pr $f\in A$, $x\in V_A(k')=X(k')$ \quad $(fg)(x)=f(x)g(x)$, \unsure{$1(x)=1_{k'}$|$(1)(x)=1_{k'}$}, $(f+g)(x)=f(x)+g(x)$}

[si $A=P_I/J$, $f\in A$ se relève en $g\in P_I$, alors \deleted{\unsure{$\longleftrightarrow$}}

\diagram{\begin{tikzcd}[ampersand replacement=\&,row sep=small]
V_A(k')\arrow[r,"f_{k'}"]\&k'\&\text{est induit par}\\
\mathbb E^I(k')=k^{\prime I}\arrow[ur,"g_{k'}"']\&\&\text{défini par }g\,]
\end{tikzcd}}

\margin{$A$ \uline{l'anneau des fonctions sur l'espace algébrique affine} $X$.}

La structure d'anneau de $A$ provient de la structure d'anneau sur le foncteur $\mathbb E_k^1$.
\scanpage{26}

\deleted{\diagram{\begin{tikzcd}[ampersand replacement=\&,row sep=large,column sep=huge]
X/F\arrow[d,phantom,"\rotatebox{90}{$\simeq$}"]\&X'/F'\arrow[d,dash]\\
\underline{\mathrm{Hom}}_X(F,\unsure{\Phi|\varphi|\mathbb L|\mathcal L}_X)\arrow[r]\&X'\\
X\arrow[u,hook]\arrow[ur]\&
\end{tikzcd}}}

\deleted{\unsure{$X$}}\quad\deleted{$\underline{\mathrm{Hom}}$}

\deleted{$F$}

\deleted{$E$ ens \unsure{ordinaire|ordinal} infini $\{\mathbb Z/\unsure{2|n}\mathbb Z\}$ \unsure{ens|ou} à deux élts}

\deleted{\unsure{$E$|$\mathscr E$} des \unsure{parties} de $E_{X'}$ en \unsure{somme|sous|sommes} de \unsure{son objet}}

\deleted{\[
\begin{aligned}
&\mathscr X/\underline{\mathrm{Hom}}(F,\unsure{G|\mathscr G|\mathbb L})\\
&X/F
\end{aligned}
\]}

\deleted{\[
\begin{aligned}
&\deleted{\unsure{\varphi|\prod}}\ \mathrm{Hom}(X',I_{X'}^E)\\
&\mathrm{Hom}(X',I_{X'}^E)
\end{aligned}
\]}

\deleted{\[
(P\amalg P)^{\deleted{\unsure{\Gamma(X')}}E}
\]}

\deleted{Topos associé à l'espace compact $I^E$}

% archive sheet number: 25
\scanpage{27}

\deleted{Sous-espace algébrique fermé d'un \unsure{($=$} espace algébrique affine.}

\textcircled{9} $\varprojlim$ dans la catégorie $\mathit{Aff}_k$ des espaces algébriques affines sur $k$.

Correspond aux $\varinjlim$ dans la catégorie $\mathit{Alg}_k$ des $k$-algèbres. Bien adapté à l'interprétation de $C$ comme sous-catégorie pleine de $\widehat C$.

$*$ Produits de schémas $\Longleftrightarrow$ ``sommes'' de $k$-algèbres $A_i$ ($i\in I$)

\[
A=\bigotimes_{i\in I}A_i
\]

\inserted{Cas particulier :} \deleted{Objet final} \quad $e_k=V_k$.

\margin{\uline{Sections} sur $k$ (ou $V_k$, ou $e_k$) $=$ pts rationnels dans \deleted{\unsure{le}} \unsure{l'interprétation} \deleted{\unsure{par}}}

$*$ Produit fibré $\Longleftrightarrow$ Produit tensoriel $B\otimes_A C$

[produit dans $C_{/S}$] \qquad (``somme dans $\mathit{Alg}_A$)

$\Bigl[\ \Updownarrow$ Changement de base $C_{/S}\to C_{/S'}$, adjoint à droite du foncteur ``oubli de $S'$ au profit de $S$''

\diagram{\begin{tikzcd}[ampersand replacement=\&,column sep=small,row sep=small]
X\arrow[d]\&X'\arrow[l]\arrow[d]\&\deleted{\unsure{Z|T}}\arrow[ll,bend right=13]\arrow[dl]\\
S\&S'\arrow[l]\&
\end{tikzcd}}

\diagram{\begin{tikzcd}[ampersand replacement=\&,column sep=small,row sep=small]
B\arrow[r]\arrow[rr,bend left=20]\&B'\arrow[r,dotted]\&C'\\
A\arrow[r]\arrow[u]\&A'\arrow[u]\arrow[ur]\&
\end{tikzcd}}

\margin{\uline{Interprétation en terme de foncteurs} ``restriction'' du foncteur des $k$-algèbres à $k'$-algèbres (ou de $A$-algèbres à $A'$-algèbres)}

Cela donne toutes les $\varprojlim$. Autres exemples

$*$ Noyaux de doubles flèches \quad $X\overset{u}{\underset{v}{\rightrightarrows}}Y$

\[
\mathrm{Ker}(u,v)=X\times_{(Y\times Y)}\Delta_Y
\]

\diagram{\begin{tikzcd}[ampersand replacement=\&,row sep=small]
X\arrow[r,"{(u,v)}"]\&Y\times Y\\
K\arrow[u]\arrow[r]\&Y\arrow[u,"\Delta_Y"']
\end{tikzcd}}

Conoyaux de double flèches \deleted{$\unsure{A}\overset u{\underset v{\rightrightarrows}}$} $B\overset u{\underset v{\leftleftarrows}}A$

\[
\simeq B/\text{idéal engendré par les } u(x)-v(x)
\]

\diagram{\begin{tikzcd}[ampersand replacement=\&,column sep=small,row sep=normal]
\bigl[\quad B\arrow[d]\&A\otimes A\arrow[l,"{(u,v)}"']\arrow[d]\&x\otimes y\arrow[d]\\
C=?\&A\arrow[l]\&xy
\end{tikzcd}}

\margin{morphisme surjectif noyau noyau engendré par $x\otimes1-1\otimes x$}

$*$ limites projectives filtrantes : correspondent aux limites inductives filtrantes de $\mathit{Alg}_k$, qui se construisent ensemblistement.

\uline{Exemple} Noyau d'un hom de groupes algébriques affines$/k$

\[
G\longrightarrow G'
\]

\uline{Ex} $\mathrm{Sl}(n)_k=\underline{\mathrm{Ker}}\bigl(\mathrm{Gl}_{\deleted{\unsure{n}}}(n)_{\unsure{k|gk}}\xrightarrow{\det}G_{\mathrm m\,k}\bigr)$. \deleted{\unsure{Ker de la dét de groupe}}

% archive sheet number: 26
\scanpage{28}

$G$ opère sur $X$, \uline{stabilisateur} d'un $x\in X(k)\unsure{=|\simeq}\mathrm{Hom}_{\mathit{Aff}_k}(e_k,X)$

\[
g\longmapsto g.x
\]

\diagram{\begin{tikzcd}[ampersand replacement=\&,row sep=small]
G\arrow[r]\&X\\
\deleted{\unsure{K}}G_{(x)}\arrow[u]\arrow[r]\&e_k\arrow[u]
\end{tikzcd}}

$\mathrm I$.

Si on veut le sous-groupe de stabilité de plusieurs éléments $x_i\in X_i(k)$, $G$ opérant sur les $X_i$, on prend $\bigcap_i G_{x_i}$ (cas particulier de produit fibré) \unsure{où|ou} on fait opérer $G$ sur $X=\prod X_i$, et on prend le stabilisateur de $x=(x_i)_{i\in I}\in X(k)$.

\textcircled{10} $\varinjlim$ dans $\mathit{Aff}_k$. Ne s'interprète pas de façon commode en terme de $\widehat{\mathit{Aff}_k}\simeq\underline{\mathrm{Hom}}(\deleted{\unsure{k}}\mathit{Alg}_k,(\mathrm{Ens}))$, et n'a pas beaucoup de sens géométriquement, sauf en des cas très particuliers.

[\uline{Exemple} $G=\mathrm{Gl}(n)_k$, $B=\mathrm{Tr}(n)_k\subset G$, alors \deleted{\unsure{B}} $B$ définit une relation d'équivalence dans $G$, et le quotient $G/B$ \uline{dans} $\mathit{Aff}_k$ est $e_k$ -- ce qui est idiot !]

% slanted
\margin{définir la relation d'équivalence associée pour un \unsure{ss-groupe|sous-groupe} d'un \unsure{gr. alg.|groupe} \ldots}

Correspond à $\varprojlim$ dans $\mathit{Alg}_k$. Ce sont les $\varprojlim$ ensemblistes (nous fait une belle jambe !)

\deleted{\unsure{a)|$*$}} \uline{Sommes dans} $\mathit{Aff}_k$. Correspond aux produits dans $\mathit{Alg}_k$. Seules les sommes finies ($\Longleftrightarrow$ produits finis) ont une interprétation \uline{géométrique simple}.

\[
X_i\longrightarrow X
\]

Attention, $\boxed{\coprod_i X_i(k')\to X(k')}$ n'est pas bijective pour $\forall\ k'$ ; \deleted{\unsure{f}} i.e.

\[
\boxed{\coprod_{\unsure{I|i}}\mathrm{Hom}_{k\text{-}\mathrm{alg}}(\deleted{\unsure{B}}A_{\unsure{i|j}},k')\longrightarrow
\mathrm{Hom}_{k\text{-}\mathrm{alg}}(A=\prod A_i,k')}
\]

n'est pas bijective. \inserted{Elle est \uline{injective} si $k'\ne0$, / \deleted{\unsure{elle n'est pas}}} Un hom de $A$ dans $k'$ ``consiste'' en la donnée d'une famille d'idempotents de $k'$ $(e'_i)_{i\in I}$ de somme $1$ et deux à deux orthogonaux et d'hom. $A_i\to k'/(1-e'_i)k'$ \unsure{?}\ldots{} Mais si

% archive sheet number: 27
\scanpage{29}

$k'$ est ``connexe'' i.e. pour tte telle décomposition de $1\in k'$, \inserted{(où $I\ne\varnothing$ et} tous les $e'_i$ sauf un sont nuls, l'autre $1$, \deleted{\unsure{donc}} (p. ex. $k'$ \deleted{\unsure{corps}} intègre, ou $k'/\mathrm{Nil}_{k'}$ intègre, \deleted{\unsure{ou plus}} ou \inserted{$k'$} un anneau local \ldots\ldots) alors les applications considérées sont bijectives.

% slanted
\margin{$*$ \uline{équivalent à la notion de connexité dans une catégorie} \ldots Exemple : $k'$ intègre}

Si $I=\varnothing$, on trouve un objet initial dans $\mathit{Aff}_k$

\[
\varnothing_k=V_0\ \text{(algèbre nulle sur $k$)}
\]

On a \deleted{\unsure{donc}}

\[
\begin{aligned}
\varnothing_k(k')&=\begin{cases}
\varnothing&\text{si }k'\ne0 \quad\text{i.e. }V_{k'}\nleftarrow\varnothing_k\\
\{e\}&\text{si }k'=0
\end{cases}\\
\rotatebox{90}{$=$}&\\
\mathrm{Hom}(V_{k'},\varnothing_k)&
\end{aligned}
\]

Donc c'est un \uline{objet initial strict}. On l'appelle le \uline{$k$-espace algébrique vide} (un \deleted{\unsure{léger|objet}} abus de langage \ldots,

$(*)$ Conoyaux $\mathrm{Coker}(u,v:R\rightrightarrows X)$ correspond à $\mathrm{Ker}(u,v:A\rightrightarrows B)$ mais n'a pas beaucoup de sens. O.K. si $R$ est une ``relation d'équivalence'' dans $X$ et si $R$ fini sur $X$ par $u$, i.e. $B$ une $A$-algèbre finie. On reviendra dessus avec la question du passage aux quotients.

Sous-espaces algébriques fermés d'un espace algébrique affine $\mathscr X\simeq V_A$ : ils correspondent aux \uline{idéaux} de $A$ avec relation d'ordre renversée

\textcircled{\unsure{11|10}} \uline{Spectre d'une $k$-algèbre $A$} : espace sous-jacent.

\[
A\longleftrightarrow\mathscr X=V_A
\]

On va associer un objet géométrique

\textcircled{a} \unsure{Elément|Ensemble} sous-jacent : ensemble \inserted{$X$} des ``lieux'' \inserted{(loci)} de ``pts géométriques'' de $\mathscr X$

\[
\Longleftrightarrow\quad\text{ensemble }\deleted{\unsure{de tous}}\inserted{\mathrm{Spec}(A)}\ \text{des idéaux premiers de } A
\]

% archive sheet number: 28
\scanpage{30}
% archive sheet number: 29

La connaissance de l'ensemble $X$ et de la famille des corps $(k(x))_{x\in X}$ (corps résiduels \deleted{\unsure{des}} en les pts de $\mathrm{Spec}\,A$ ou $\mathrm{Spec}(\mathscr X)$ -- $k(x)$ représente le foncteur $k'$ (corps variable sur $k$) $\longmapsto$ ens des \deleted{points} éléments de $\mathscr X(k')$ ayant pour lieu $x$).

\textcircled{b} \uline{Topologie sur $X$} (les fermés sont les) $V(S)\subset X$ associés aux parties $S\inserted{=(f_\alpha)}$ de $A$, \deleted{\unsure{leur}} un point \inserted{géom. $u$} à valeurs dans $k'$ a son lieu dans $V(S)$ sss $f_\alpha(u)=0$ pour tout $\alpha$. On a \deleted{$V(S)=V(J_S)$ si} $V(S)=$ ensemble des idéaux premiers de $A$ qui contiennent $S$.

On a $V(S)=V(J_S)$, $J_S$ idéal engendré, et $\unsure{\overline{\mathrm m}|=}$ $V(S)=V(\widetilde J_S)$. Il \deleted{\unsure{On a $V($|Donc}} est bon de se borner aux idéaux, et $\overline{\mathrm m}$ aux idéaux ``radicaux''. \deleted{L'ensemble des} On voit que $V(\unsure{S|J})=V(\unsure{S'|J'})$ (resp. $V(S)\subset V(S')$) signifie, avec des notations précédentes, que $V'_S=V^{\unsure{\prime\prime|\prime}}_{S'}$ (restrictions aux corps) \deleted{du moins que} (resp $V'_S\deleted{\unsure{(k)}}\subset V'_{S'}$), donc $\widetilde J_S=\widetilde J_{S'}$ (resp. $\widetilde J_S\supset\widetilde J_{S'}$) : cela provient au fond que $\widetilde J_S$ est l'intersection des idéaux premiers qui contiennent $A$.

\[
\deleted{V(\sum J_i}\ V(\bigcup_{i\in I}S_i)=V(\sum J_i)=\bigcap_i V(S_i)=\bigcap V(J_i)
\qquad I\ \text{quelc.}
\]

(NB une \deleted{intersection} somme de deux \inserted{idéaux} \deleted{\unsure{rad}} radicaux n'est pas nécessairement un idéal radical)

\[
V(\bigcap_i J_i)=V(\prod J_i)=\bigcup_i V(J_i)\qquad I\ \text{fini}
\]

[NB $\prod_i J_i\subset\bigcap_i J_i$, il faut prouver que $\forall$ idéal premier $\mathfrak p$ qui contient l'\deleted{intersection}\inserted{produit} des $J_i$ contient un des $J_i$, or sinon on aurait pour $\forall\ i\in I$ un $f_i\in J_i,\notin\mathfrak p$, \deleted{on en passant à $A/\mathfrak p$ que si un produit des idéaux} donc $\prod f_i$ serait $\in\bigcap J_i$ et $\notin\mathfrak p$, absurde ] \deleted{\unsure{premiers est nul}, l'un d'eux est nul}

Donc on a une \uline{topologie}.

\margin{Si $x\in X$ correspond à $p_x\subset A$, Alors $\overline{\{x\}}$ est le plus petit fermé $V(J)$ qui contient $x$, i.e. tel que $J\subset p_x$, donc \deleted{\unsure{plus}} c'est $V(p_x)$.}

\uline{Espace \deleted{\unsure{top.}} quasi-compact} \inserted{(mais non séparé)}

\margin{$x\in\overline{\{y\}}$ $\Updownarrow$ $\overline{\{x\}}\subset\overline{\{y\}}\Longleftrightarrow p_x\supset p_y$}

$[\,X_f=X-V(f)$ est ouvert, \inserted{(le lieu de $u\in\mathscr X(k)$ est dans $X_f$ sss $f(u)\ne0$)}, et les \deleted{\unsure{$D(f)$}} $X_f$ forment une base de la topologie de $X$, stable par intersections finies ($X_{1}=X$, $X_{fg}=X_f\cap X_g$).

% slanted
\margin{plus bas}
\scanpage{31}

\[
\bigl[\,X\ \text{quasi-compact}\Longleftrightarrow
\text{Si } A=\sum_{i\in I}J_i,\ \exists\,J\ \text{fini}\subset I,\ A=\sum_{i\in J}J_i.
\]

\deleted{\unsure{Parties irréductibles}}

Si $J$ est un idéal de $A$, $\mathrm{Spec}\,(A/J)=Y\subset X$ est un sous-espace topologique fermé de $X$. \deleted{On trouve} Il ne change pas si on remplace $J$ par $\widetilde J$, donc \deleted{\unsure{pour étu}} un sous-espace fermé de $X$ est tjrs hom. au spectre d'un anneau réduit.

\uline{Prop} Si $A$ est réduit, $X=\mathrm{Spec}\,A$ irréductible $\Longleftrightarrow A$ intègre.

\[
\begin{aligned}
&\text{En effet, $\mathrm{Spec}\,A$ irréductible}\ \overset{\text{déf}}{\Longleftrightarrow}
\left\{\begin{aligned}
&X\ne\varnothing\\
&\underset{X'\ \mathrm{et}\ X''\ \mathrm{fermés}}{X=X'\cup X''}\Rightarrow X'\text{ ou }X''=X
\end{aligned}\right.\\[-1pt]
&\qquad\begin{array}{l|l}
\text{i.e.}& A\ne0\\
&\text{i.e. }J'\cdot J''=0\Rightarrow J'\text{ ou }J''\ \text{nul}
\end{array}
\end{aligned}
\]

\uline{Corollaire} Soit $A$ quelconque, \deleted{\unsure{il y a correspondance biun|$A$ a une correspondance biunivoque}} les parties fermées irréductibles de $A$ sont les $V(p)$, $p$ idéal premier de $A$ (correspondance biunivoque \unsure{et|$\updownarrow$} renversant l'ordre) et \deleted{\unsure{l'id}} si $x$ correspond à $p$, alors $V(p)$ est égal à $\overline{\{x\}}$. En particulier, $X$ est \uline{sobre}.

\margin{$\mathrm{Spec}(A_f)\simeq X_f$ (homéomorphisme)}

% slanted
\margin{\textcircled{$10'$} Fonctorialité de \unsure{l'esp.} top. $\mathrm{Spec}\,A$ Cas des épi. d'ann.}

\uline{NB} L'espace topologique $\mathrm{Spec}\,A$ ne dépend que de la structure d'anneau de $A$, i.e. de $\mathbb Z$-algèbre \ldots.

\textcircled{11} \uline{Spectre d'une $k$-algèbre : faisceau d'anneaux}

On ne peut réconstituer $A$ à partir du seul espace topologique $X=\mathrm{Spec}\,A$ !

Faisceaux sur $X$ : $\forall f\in A$, ensemble $P(f)$

$\forall f,g\in A$ tels que $\deleted{\unsure{D}}X_{\unsure{fS|f}}\subset \deleted{\unsure{D}}X_g)$, une application $P(g)\to P(f)$

\[
\begin{aligned}
&\text{i.e. }V(f)\supset V(g)\\
&\text{i.e. }\widetilde{fA}\subset\widetilde{gA}\\
&\text{i.e. }f\in\widetilde{gA}\\
&\text{i.e. }\left\{\begin{aligned}&\exists\,n\in\mathbb N,\ f^n=gh\\&h\in A\end{aligned}\right.\ \big|\ \text{écrivons } f\prec g
\end{aligned}
\]

avec transitivité si $f,g,h$ tels que $f\prec g\prec h$ et identité si $f=g$

Avec $\forall f\in A$ et $(f_i)_{i\in I}$, $f_i\in{I}$ tels que

\[
\begin{aligned}
&X_f=\bigcup X_{f_i}\\
&\text{i.e. }V(f)=\bigcap V(f_i)\\
&\text{i.e. }\widetilde{fA}=\widetilde{\sum f_i A}\\
&\text{i.e. }\left\{\begin{aligned}
&\forall i\quad f_i\prec f\quad\unsure{\vee|\checkmark}\\
&\exists n\in\mathbb N\text{ et }(h_i)_{i\in I}\text{ avec } f^n=\sum f_i h_i
\end{aligned}\right.
\end{aligned}
\]

on veut exactitude dans

\[
P(f)\longrightarrow\prod_i P(f_i)\rightrightarrows\prod_{ij}P(f_i f_j)
\]

% archive sheet number: 30
\scanpage{32}

[\uline{NB} On peut supposer $I$ fini si on y tient ]

Soit $M$ un $A$-module. On définit, pour $f\in A$

\[
M_f=\varinjlim_n(M,n)\qquad\text{où }\underset{\substack{\Vert\\M}}{(M,n)}
\xrightarrow{f^{n'-n}}\underset{\substack{\Vert\\M}}{(M,n')}\quad\text{si }n'\geq n
\]

Donc les éléments de $M_f$ sont les $m/f^n$, avec règles de calcul habituelles. C'est la solution au pb universel d'envoyer $M$ dans un $A$-module où $f$ opère de façon inversible.

\[
M\xrightarrow{i_f}M_f
\]

\unsure{Le} \uline{Lemme} Si $f\prec g$, $\exists$ un unique hom de $M_g$ dans $M_f$ compatible avec $i_g,i_f$

\diagram{\begin{tikzcd}[ampersand replacement=\&,row sep=small,column sep=tiny]
\&M\arrow[dl,"i_g"']\arrow[dr,"i_f"]\&\\
M_g\arrow[rr,"{i_{f,g}}"]\&\&M_f
\end{tikzcd}}

En effet, $f$ devenant inversible dans $M_f$ donc $f^n=gh$ aussi, $g$ devient inversible dans $M_f$, et c'est \unsure{là|cela} l'assertion.

% slanted
\margin{\uline{NB} $A_f$ une $A$-algèbre \unsure{à lte}, $A_f\to A_g$ un \unsure{hom d'algèbres|hom de} \quad $M_f$ un $A_f$-module [$A_f$ \deleted{\unsure{dualité}} $\simeq A_f$ \unsure{pr.}!] $M_f\simeq M\otimes_A A_f$}

\uline{Cor} Transitivité : on a un préfaisceau \inserted{$\widetilde M$} sur \deleted{la base des} catégorie \deleted{rigide} raide $A^{\unsure{\bullet|\circ}}$. En fait, c'est un préfaisceau de modules sur le préfaisceau d'anneaux $\unsure{f|g}\mapsto A_f$.

\uline{Théorème} Le préfaisceau précédent est un \uline{faisceau}.

\[
\begin{aligned}
M_f\longrightarrow\prod M_{f_i}\rightrightarrows\prod M_{f_i f_j}
&\qquad\text{exact si}\\
& f_i\prec f\text{ et}\\
& f^n=\sum f_i h_i
\end{aligned}
\]

% slanted
\margin{Si $g\prec f$, on a $M_g=(M_f)_g$ transitivité. \unsure{des}}

On est ramené au cas où \deleted{\unsure{$f=1$, $A$, $M_f$}} \inserted{$f=1$} (car \deleted{\unsure{Al}} $f\succ g$ implique $(M_f)_g\simeq M_g$) \ldots.

En particulier, $X$ est \uuline{annelé par} $\widetilde A$, noté $\unsure{\underline O|\mathcal O|\underline{\mathcal O}}_X$, faisceau \deleted{\unsure{d'algèbres}} d'anneaux (et $=$ de $A$-algèbres).

Fibres du faisceau $M_{\mathfrak p}$, $A_{\mathfrak p}=$

\[
A_{\mathfrak p}=\underline O_{X,x}\quad\text{\uline{anneau local}}
\]

\margin{[$S\subset A$ stable par multiplication donc filtrant pour \deleted{\unsure{\prec}} $\succ$ donc $MS^{-1}=\varinjlim_{f\in S}M_f=\{\frac{x}{f^n}$ \unsure{avec} règles de calcul \ldots.]}

% archive sheet number: 31
